%==============================================================================
% Major Project - Interim Progress Report
% Title : SON-GNN-DQN: A Topology-Generalized Graph Neural Network and
%         Deep Q-Learning Framework for Predictive Handover and Load
%         Optimization in LTE Networks
% Authors: Rojin Puri, Rujan Subedi, Saange Tamang, Sulav Kandel
% Dept.  : Electronics, Communication and Information Engineering
% Campus : Paschimanchal Campus, IOE, Tribhuvan University
%
% NOTE: This document is the interim progress report submitted to the
% Department of Electronics and Computer Engineering and shared with
% Nepal Telecom (NTC) Pokhara as a basis for requesting drive-test and
% OSS data for the final validation phase. All results reported here are
% from the synthetic simulator that the team has built and verified;
% drive-test and NTC validation results will be added in the final
% submission.
%==============================================================================

\documentclass[12pt,a4paper,oneside]{report}

%------------------------------- PACKAGES -------------------------------------
\usepackage[utf8]{inputenc}
\usepackage[T1]{fontenc}
\usepackage{lmodern}
\usepackage{times}
\usepackage[a4paper,left=1.25in,right=1in,top=1in,bottom=1in]{geometry}
\usepackage{setspace}
\onehalfspacing
\usepackage{graphicx}
\usepackage{float}
\usepackage{subcaption}
\usepackage{amsmath,amssymb,amsfonts}
\usepackage{bm}
\usepackage{booktabs}
\usepackage{longtable}
\usepackage{multirow}
\usepackage{array}
\usepackage{tabularx}
\usepackage{caption}
\captionsetup{font=small,labelfont=bf,justification=centering}
\usepackage{enumitem}
\usepackage{hyperref}
\hypersetup{colorlinks=false,pdfborder={0 0 0}}
\usepackage{xcolor}
\usepackage{listings}
\usepackage{titlesec}
\usepackage{fancyhdr}
\usepackage{tocloft}
\usepackage{algorithm}
\usepackage{algpseudocode}
\usepackage{pgfplots}
\pgfplotsset{compat=1.17}
\usepackage{tikz}
\usetikzlibrary{positioning,shapes,arrows.meta,calc,fit,backgrounds}

%------------------------------- LISTINGS STYLE -------------------------------
\definecolor{codegray}{rgb}{0.95,0.95,0.95}
\definecolor{codeblue}{rgb}{0.13,0.29,0.53}
\definecolor{codegreen}{rgb}{0.0,0.5,0.0}
\lstset{
  backgroundcolor=\color{codegray},
  basicstyle=\ttfamily\footnotesize,
  keywordstyle=\color{codeblue}\bfseries,
  commentstyle=\color{codegreen}\itshape,
  stringstyle=\color{red!60!black},
  numbers=left,
  numberstyle=\tiny\color{gray},
  stepnumber=1,numbersep=8pt,
  showspaces=false,showstringspaces=false,
  frame=single,rulecolor=\color{black!30},
  breaklines=true,captionpos=b,
  language=Python
}

%------------------------------- TITLE FORMATTING ------------------------------
\titleformat{\chapter}[block]
  {\normalfont\Large\bfseries}
  {\thechapter.}{0.8em}{\Large\MakeUppercase}
\titlespacing*{\chapter}{0pt}{-20pt}{20pt}

\titleformat{\section}
  {\normalfont\large\bfseries}{\thesection.}{0.6em}{}
\titleformat{\subsection}
  {\normalfont\normalsize\bfseries}{\thesubsection.}{0.5em}{}

\pagestyle{plain}

\renewcommand{\cfttoctitlefont}{\Large\bfseries\MakeUppercase}
\renewcommand{\cftloftitlefont}{\Large\bfseries\MakeUppercase}
\renewcommand{\cftlottitlefont}{\Large\bfseries\MakeUppercase}
\renewcommand{\cftchapfont}{\bfseries}
\renewcommand{\cftchappagefont}{\bfseries}
\setlength{\cftbeforechapskip}{6pt}

%==============================================================================
\begin{document}
%==============================================================================

%------------------------------ TITLE PAGE -------------------------------------
\begin{titlepage}
\begin{center}
    \vspace*{0.5cm}
    \includegraphics[width=0.22\textwidth]{figures/tu_logo.png}\\[0.8cm]
    {\large\bfseries TRIBHUVAN UNIVERSITY}\\[0.15cm]
    {\large\bfseries INSTITUTE OF ENGINEERING}\\[0.15cm]
    {\large\bfseries PASCHIMANCHAL CAMPUS}\\[1.8cm]

    {\large\bfseries A Major Project Progress Report}\\[0.25cm]
    {\large\bfseries On}\\[0.4cm]
    {\large\bfseries SON-GNN-DQN: A Topology-Generalized Graph Neural
    Network and Deep\\ Q-Learning Framework for Predictive Handover and
    Load Optimization\\ in LTE Networks}\\[2.0cm]

    {\bfseries Submitted By:}\\[0.4cm]
    \begin{tabular}{ll}
        Rojin Puri      & PAS078BEI028 \\[2pt]
        Rujan Subedi    & PAS078BEI030 \\[2pt]
        Saange Tamang   & PAS078BEI031 \\[2pt]
        Sulav Kandel    & PAS078BEI041 \\
    \end{tabular}\\[1.8cm]

    {\bfseries Submitted To:}\\[0.3cm]
    Department of Electronics and Computer Engineering\\
    Paschimanchal Campus\\
    Lamachour, Pokhara--16, Nepal\\[1.5cm]

    {\bfseries May, 2026}
\end{center}
\end{titlepage}

%---------------------- FRONT MATTER (roman numerals) --------------------------
\pagenumbering{roman}
\setcounter{page}{2}

%------------------------------ DECLARATION ------------------------------------
\chapter*{DECLARATION}
\addcontentsline{toc}{chapter}{DECLARATION}

We, the undersigned students of the Bachelor of Electronics, Communication
and Information Engineering at Paschimanchal Campus, Institute of
Engineering, Tribhuvan University, hereby declare that the project entitled
\textit{``SON-GNN-DQN: A Topology-Generalized Graph Neural Network and Deep
Q-Learning Framework for Predictive Handover and Load Optimization in LTE
Networks''} is our original work carried out under the supervision of
Asst.\ Prof.\ Er.\ Hom Nath Tiwari, Head of the Department of Electronics
and Computer Engineering, Paschimanchal Campus.

This is an interim progress report. The work presented here has been
implemented and verified through synthetic simulation. The final
validation phase, which depends on real drive-test and operator-side
measurements, is in progress. All sources of information used in this
report have been duly acknowledged. The document has not been submitted,
in part or in full, to any other university or institution for the award
of any degree.

\vspace{1.5cm}

\begin{tabular}{p{6cm}p{6cm}}
Rojin Puri \dotfill   & Rujan Subedi \dotfill \\[1.0cm]
(PAS078BEI028)        & (PAS078BEI030) \\[1.5cm]
Saange Tamang \dotfill & Sulav Kandel \dotfill \\[1.0cm]
(PAS078BEI031)        & (PAS078BEI041) \\
\end{tabular}

\vspace{1cm}
\noindent Date: May, 2026

%------------------------------ CERTIFICATE ------------------------------------
\chapter*{CERTIFICATE OF APPROVAL}
\addcontentsline{toc}{chapter}{CERTIFICATE OF APPROVAL}

The undersigned certify that they have read and recommended this interim
project progress report entitled \textit{``SON-GNN-DQN: A Topology-
Generalized Graph Neural Network and Deep Q-Learning Framework for
Predictive Handover and Load Optimization in LTE Networks''}, submitted by
Rojin Puri, Rujan Subedi, Saange Tamang and Sulav Kandel, in partial
fulfilment of the requirements for the Bachelor's degree in Electronics,
Communication and Information Engineering.

\vspace{2cm}
\noindent\rule{6cm}{0.5pt}\\
\textbf{Project Supervisor}\\
Asst.\ Prof.\ Er.\ Hom Nath Tiwari\\
Head of Department\\
Department of Electronics and Computer Engineering\\
Paschimanchal Campus, IOE

\vspace{1.5cm}
\noindent\rule{6cm}{0.5pt}\\
\textbf{External Examiner}\\
\textit{[Name and Designation]}

\vspace{1.5cm}
\noindent\rule{6cm}{0.5pt}\\
\textbf{Internal Examiner}\\
\textit{[Name and Designation]}

\vspace{1cm}
\noindent Date: May, 2026

%------------------------------ COPYRIGHT --------------------------------------
\chapter*{COPYRIGHT}
\addcontentsline{toc}{chapter}{COPYRIGHT}

The authors have agreed that the library of the Department of Electronics
and Computer Engineering, Paschimanchal Campus, Institute of Engineering,
may make this report freely available for inspection. The authors have
further agreed that permission for extensive copying of this report for
scholarly purposes may be granted by the supervisor who supervised the
project, or, in his absence, by the Head of the Department. It is
understood that recognition will be given to the authors and to the
Department of Electronics and Computer Engineering, Paschimanchal Campus,
in any use of the material. Copying, publication or other use of this
report for financial gain without the written approval of the Department
and the authors is strictly prohibited.

Requests for permission to copy or otherwise use the material in this
report, in whole or in part, should be addressed to:

\vspace{1cm}
\noindent\hspace{1cm}\textit{Head}\\
\noindent\hspace{1cm}\textit{Department of Electronics and Computer
Engineering}\\
\noindent\hspace{1cm}\textit{Paschimanchal Campus, Institute of Engineering}\\
\noindent\hspace{1cm}\textit{Lamachour, Pokhara--16, Nepal}

%---------------------------- ACKNOWLEDGEMENT ----------------------------------
\chapter*{ACKNOWLEDGEMENT}
\addcontentsline{toc}{chapter}{ACKNOWLEDGEMENT}

We are pleased to present this interim progress report on the project
titled \textit{``SON-GNN-DQN: A Topology-Generalized Graph Neural Network
and Deep Q-Learning Framework for Predictive Handover and Load
Optimization in LTE Networks''}. We extend our sincere gratitude to the
Department of Electronics and Computer Engineering, Institute of
Engineering, Paschimanchal Campus, for providing us with the opportunity
to undertake this major project as part of our academic curriculum.

We would like to express our deepest appreciation to our project
supervisor, Asst.\ Prof.\ Er.\ Hom Nath Tiwari, Head of the Department of
Electronics and Computer Engineering, for his invaluable guidance,
constant motivation, and thoughtful feedback throughout every stage of
this project. His insights into mobility management, reinforcement
learning, and academic writing have significantly shaped the technical
direction and rigor of this work.

We are also thankful to the faculty members of the Department of
Electronics and Computer Engineering, Paschimanchal Campus, for their
continuous encouragement and academic support. We wish to extend our
appreciation in advance to the engineers of Nepal Telecom (NTC), Pokhara
Branch, with whom we hope to collaborate during the validation phase of
this project by sharing drive-test logs and operator-side handover
parameters for academic, non-commercial use.

We are indebted to our families, friends, and well-wishers whose
encouragement has been a constant source of strength throughout this
academic journey.

\vspace{1.2cm}
\begin{tabular}{ll}
Rojin Puri      & PAS078BEI028 \\[5pt]
Rujan Subedi    & PAS078BEI030 \\[5pt]
Saange Tamang   & PAS078BEI031 \\[5pt]
Sulav Kandel    & PAS078BEI041 \\
\end{tabular}

%------------------------------ ABSTRACT ---------------------------------------
\chapter*{ABSTRACT}
\addcontentsline{toc}{chapter}{ABSTRACT}

Handover management in Long-Term Evolution (LTE) networks continues to
rely on threshold-based, rule-driven schemes such as the Event~A3 trigger
with a fixed Time-to-Trigger (TTT) and a static Cell Individual Offset
(CIO). Although standardized by 3GPP, these mechanisms struggle to adapt
to dense heterogeneous deployments, hilly terrain, and high mobility
scenarios that are characteristic of Nepalese cities such as Pokhara.
This work proposes \textbf{SON-GNN-DQN}, a topology-generalized Self-
Organizing Network (SON) compatible framework that combines a Graph
Neural Network (GNN) preference model with a Deep Q-Learning (DQN)
decision agent and a safety-bounded SON translation layer that converts
learned preferences into bounded CIO and TTT adjustments. The GNN encodes
the cellular topology as a dynamic graph in which base stations and the
User Equipment (UE) form nodes and neighbor relations form edges. The
DQN selects target cells from the learned preference embedding. A
\texttt{ue\_only} feature profile, based on measurements that a single
device can collect, is adopted as the main path so that the framework
can later be validated with single-UE drive-test data and an extended
\texttt{oran\_e2} profile is retained for future operator-side
integration.

This interim report documents the system design, the implemented
simulator, the GNN-DQN agent with replay buffer and SON translation
layer, the baseline policies, the unit-test suite, and the short
training runs (smoke and diagnostic) executed on the synthetic
environment. The long multi-scenario training run was started on the
available hardware (Apple~M1 Pro, 16~GB unified memory) and Google Colab
Pro; results from this run will be reported in the final submission.
Drive-test validation in Pokhara and operator-side parameter validation
with Nepal Telecom are planned for the next phase of the project, and a
formal data-request letter is included as Appendix~D.

\vspace{0.6cm}
\noindent\textbf{Keywords:} LTE Handover, Graph Neural Network, Deep
Q-Learning, Self-Organizing Network, Cell Individual Offset, Mobility
Load Balancing, Drive Test, Pokhara, Nepal Telecom.

%------------------------------ TOC, LOF, LOT ----------------------------------
\renewcommand{\contentsname}{TABLE OF CONTENTS}
\tableofcontents
\clearpage

\renewcommand{\listfigurename}{LIST OF FIGURES}
\listoffigures
\clearpage

\renewcommand{\listtablename}{LIST OF TABLES}
\listoftables
\clearpage

%---------------------------- LIST OF ABBREVIATIONS ----------------------------
\chapter*{LIST OF ABBREVIATIONS}
\addcontentsline{toc}{chapter}{LIST OF ABBREVIATIONS}

\begin{center}
\begin{tabular}{ll}
\toprule
\textbf{Abbreviation} & \textbf{Full Form}\\
\midrule
3GPP        & 3rd Generation Partnership Project\\
A3          & Event A3 (Neighbor becomes better than serving)\\
BS          & Base Station\\
CIO         & Cell Individual Offset\\
CPU         & Central Processing Unit\\
CSV         & Comma Separated Values\\
DL          & Deep Learning\\
DQN         & Deep Q-Network\\
DRL         & Deep Reinforcement Learning\\
EARFCN      & E-UTRA Absolute Radio Frequency Channel Number\\
eNB         & Evolved Node B (4G LTE Base Station)\\
EPC         & Evolved Packet Core\\
E2          & Interface between RAN and Near-RT RIC in O-RAN\\
GAT         & Graph Attention Network\\
GCN         & Graph Convolutional Network\\
gNB         & Next Generation Node B (5G Base Station)\\
GNN         & Graph Neural Network\\
GPU         & Graphics Processing Unit\\
HO          & Handover\\
KPI         & Key Performance Indicator\\
KPM         & Key Performance Measurement\\
LTE         & Long-Term Evolution\\
MAC         & Medium Access Control\\
MLB         & Mobility Load Balancing\\
ML          & Machine Learning\\
MRO         & Mobility Robustness Optimization\\
NTC         & Nepal Telecom\\
O-RAN       & Open Radio Access Network\\
OSS         & Operations Support System\\
PCI         & Physical Cell Identifier\\
PRB         & Physical Resource Block\\
PyG         & PyTorch Geometric\\
QoS         & Quality of Service\\
RAN         & Radio Access Network\\
RIC         & RAN Intelligent Controller\\
RL          & Reinforcement Learning\\
RLF         & Radio Link Failure\\
RRC         & Radio Resource Control\\
RSRP        & Reference Signal Received Power\\
RSRQ        & Reference Signal Received Quality\\
SINR        & Signal-to-Interference-plus-Noise Ratio\\
SON         & Self-Organizing Network\\
TTT         & Time-to-Trigger\\
UE          & User Equipment\\
xApp        & Application running on Near-RT RIC in O-RAN architecture\\
\bottomrule
\end{tabular}
\end{center}

%---------------------- MAIN MATTER (arabic numerals) --------------------------
\clearpage
\pagenumbering{arabic}
\setcounter{page}{1}

%==============================================================================
\chapter{INTRODUCTION}
%==============================================================================

\section{Background}

The telecommunications sector of Nepal has expanded rapidly over the past
decade. Long-Term Evolution (LTE) is now the primary technology for
nationwide mobile broadband, deployed by Nepal Telecom (NTC) and Ncell in
Band~3 (1800~MHz) and Band~1 (2100~MHz) across urban centres, highways
and selected hilly clusters. Although the volume of LTE deployment has
grown, the mobility-management procedures used in the field remain
largely tied to the 3GPP Release~8/9 baseline: a Reference Signal
Received Power (RSRP) measurement, an Event~A3 trigger, a fixed Time-to-
Trigger (TTT) typically in the $160$--$320$~ms range, and a globally
tuned Cell Individual Offset (CIO).

These conventional, threshold-driven handover procedures perform
reasonably well in flat, lightly loaded scenarios. However, in the
mixed urban / highway / hilly geography around Pokhara they are known to
produce ping-pong handovers when RSRP values of adjacent cells cross
repeatedly, late handovers when a UE clings to a degrading cell until a
Radio Link Failure (RLF) occurs, and load imbalance when a strong macro
cell becomes congested while neighbouring cells remain lightly loaded
\cite{3gpp_36839,zhang_mlb_2021}.

The limitations of static threshold tuning have motivated a global
research push toward AI-driven Self-Organizing Networks (SON) and toward
learning-based xApps deployed on the Near-Real-Time RAN Intelligent
Controller (Near-RT RIC) of the Open Radio Access Network (O-RAN)
architecture \cite{oran_alliance_wg2,polese_oran_survey}. Graph Neural
Networks (GNNs) have emerged as a natural representation for cellular
networks because base stations form a graph with neighbour relations as
edges and message passing captures interference and load coupling that
vector-based models cannot \cite{kipf_gcn,velickovic_gat,hamilton_graphsage}.
Deep Q-Learning (DQN) provides a model-free reinforcement learning
framework that suits sequential mobility decisions where each handover
affects future RSRP, SINR and throughput trajectories
\cite{mnih_dqn_2015}. The combination has been explored for resource
allocation \cite{wang_gnn_rl_2022}, vehicular routing
\cite{he_vanet_gnn} and, recently, handover optimization in 5G networks
\cite{tang_handover_drl,zhao_gnn_handover}. The published works typically
rely on simulators (for example, ns--3 with the LENA module) or full
O-RAN testbeds and are rarely validated with field measurements from a
developing-country LTE deployment.

This project develops SON-GNN-DQN, a topology-generalized GNN-DQN
framework tailored to the constraints faced by a small academic team in
Nepal. The framework treats the GNN as a preference model over candidate
target cells and the DQN as the temporal decision agent. A \emph{safety-
bounded SON translation layer} converts the agent's preferences into
bounded CIO and TTT adjustments that comply with the 3GPP CIO range and
operator-defined TTT values. Because the design is built around UE-only
measurements (RSRP, RSRQ, SINR, speed, serving-cell identifier), it can
be evaluated in a later phase using a single-UE drive-test dataset.

\section{Motivation}

The motivation for this work originates from three observations made
during the proposal phase. First, residents of Pokhara routinely
experience call drops, throughput variation on moving vehicles and
inconsistent video calls along the Lakeside--Mahendrapul corridor.
Although individually small, these issues aggregate into a measurable
degradation of digital experience for a city that hosts a large student
and tourist population. Second, public discussions with academic peers
and informal exchanges with telecommunications engineers in Pokhara
indicate that AI- or ML-based mobility optimization is not currently in
production. Third, recent advances in Graph Neural Networks and Deep
Reinforcement Learning have matured to a point where a credible
implementation is possible on a personal laptop and on Google Colab Pro
GPUs, making the topic feasible for a final-year engineering project at
a public university.

\section{Problem Statement}

Despite a decade of LTE deployment in Nepal and the maturity of 3GPP
mobility procedures, handover performance in cities such as Pokhara
continues to suffer from ping-pong events, late handovers, radio link
failures and load imbalance, especially in dense urban clusters, along
highway corridors and in hilly shadow-prone areas. Conventional rule-
based mechanisms -- Event~A3 with a fixed Time-to-Trigger and a globally
tuned Cell Individual Offset -- do not exploit the spatial topology of
the network or the temporal dynamics of UE mobility. There is therefore
a need for an intelligent, topology-aware and load-aware handover policy
that operates within the constraints of current Self-Organizing Network
interfaces, that can be validated using UE-side measurements without
operator OSS access, and that offers a practical evolutionary path
toward future O-RAN-based deployment.

\section{Objective}

To design, implement and validate a topology-generalized Graph Neural
Network and Deep Q-Learning (GNN-DQN) framework with a safety-bounded
SON translation layer for adaptive A3 and Cell Individual Offset (CIO)
tuning in LTE networks, and to compare its performance against
conventional handover baselines using a synthetic LTE simulator and, in
the final phase, real drive-test data collected in Pokhara.

\section{Scope and Limitations}

The scope of this interim phase covers (i) the design and implementation
of an LTE mobility simulator that produces RSRP, RSRQ, SINR, throughput,
PRB load and handover events on configurable topologies; (ii) the
construction of dynamic graph representations from UE-side measurements;
(iii) the implementation of the GNN-DQN agent under the
\texttt{ue\_only} feature profile; (iv) a safety-bounded SON translation
layer that maps learned preferences to bounded CIO and TTT adjustments;
and (v) baseline implementations (\texttt{no\_handover},
\texttt{random\_valid}, \texttt{strongest\_rsrp}, \texttt{a3\_ttt},
\texttt{load\_aware}, \texttt{gnn\_dqn}, \texttt{son\_gnn\_dqn}).

The work is bound by the following limitations. A field drive-test
campaign has not yet been carried out. A small pilot logging session was
performed on a Poco~X3 (Android) with Network Cell Info Lite to verify
that the unified CSV schema and the data-ingestion pipeline work end-to-
end, but the volume of collected data is not sufficient for statistical
evaluation; a representative drive-test campaign is planned for the
validation phase. Per-cell Physical Resource Block (PRB) utilization is
not directly observable from a UE and is represented in the simulator
by an explicit Poisson-driven load process; in the final phase this
quantity will be sourced either from NTC operator data or from a UE-side
RSRQ-based load proxy. The current results are therefore preliminary and
based on the synthetic simulator only.

%==============================================================================
\chapter{LITERATURE REVIEW}
%==============================================================================

Mobility management in cellular networks has matured along two parallel
tracks: a standardization track driven by 3GPP and the O-RAN Alliance,
and an academic track driven by the rise of Deep Reinforcement Learning
(DRL) and Graph Neural Networks (GNN). This chapter reviews both tracks
to position SON-GNN-DQN within the existing body of work.

\section{Threshold-Based Handover in 3GPP LTE}

3GPP TS~36.331 \cite{ts36331} defines the LTE Radio Resource Control
procedures and specifies the measurement events A1--A5 and B1--B2. The
most widely deployed trigger is Event~A3, in which the UE reports a
measurement when the offset-corrected RSRP of a neighbour exceeds that
of the serving cell by a margin $A3_{\text{offset}}$ for a duration
TTT. A companion technical report, 3GPP TR~36.839 \cite{3gpp_36839},
analyses mobility performance under heterogeneous deployments and
identifies ping-pong and late handover as the dominant failure modes.
Although the standard supports per-cell offsets and hysteresis, operator
practice tends toward globally tuned values, reducing adaptivity.
Castro-Hernandez and Paranthaman \cite{castro_handover_2014} and
Lopez-Perez \textit{et~al.}~\cite{lopezperez_mobility_2012} provide
measurement studies showing that the optimal
$(A3_{\text{offset}}, \text{TTT})$ pair depends on UE speed, traffic
load and inter-site distance.

\section{Self-Organizing Networks and Mobility Load Balancing}

3GPP TS~36.902 \cite{ts36902} and a series of industry white papers
introduced the SON framework with Mobility Robustness Optimization (MRO)
and Mobility Load Balancing (MLB) as core use cases. Zhang
\textit{et~al.}~\cite{zhang_mlb_2021} survey MLB algorithms and observe
that classical solutions adjust CIO based on simple PRB utilization
comparisons, which oscillate under bursty traffic. Hasan
\textit{et~al.}~\cite{hasan_son_2020} propose a fuzzy controller to
combine MRO and MLB but report difficulty in hand-tuning fuzzy rules
across heterogeneous clusters.

\section{Reinforcement Learning for Handover}

Mnih \textit{et~al.}~\cite{mnih_dqn_2015} introduced DQN, demonstrating
the viability of model-free deep RL in high-dimensional state spaces.
Several recent works apply RL to handover. Tang \textit{et~al.}
\cite{tang_handover_drl} use a DQN to learn a per-UE handover policy in
a small-cell LTE-A deployment. Wang \textit{et~al.}
\cite{wang_handover_drl_2022} combine policy-gradient methods with
long-short-term memory (LSTM) state encoders to handle fast fading.
These works rely on flat state vectors and do not exploit the network
topology.

\section{Graph Neural Networks for Wireless Networks}

GNNs are now a standard tool for representation learning on non-
Euclidean structures. Kipf and Welling \cite{kipf_gcn} introduced the
Graph Convolutional Network (GCN); Veli\v{c}kovi\'c
\textit{et~al.}~\cite{velickovic_gat} extended it with attention; and
Hamilton \textit{et~al.}~\cite{hamilton_graphsage} demonstrated
inductive generalization to unseen graphs through GraphSAGE. In the
wireless domain, Eisen and Ribeiro \cite{eisen_random_graphs_2020}
showed that GNNs are permutation-equivariant and naturally suited to
topology-varying problems. Shen
\textit{et~al.}~\cite{shen_gnn_resource_2021} applied GNNs to power
control and link scheduling; Zhao
\textit{et~al.}~\cite{zhao_gnn_handover} use GNN embeddings as the
input to a value network for handover in dense 5G networks.

\section{O-RAN, Near-RT RIC and xApps}

The O-RAN Alliance has standardized a disaggregated RAN architecture
with the Near-RT RIC as the host for xApps that consume KPM reports and
emit RAN Control (RC) actions over the E2 interface
\cite{oran_alliance_wg2}. Polese
\textit{et~al.}~\cite{polese_oran_survey} provide a comprehensive
survey of AI/ML xApps and explicitly identify mobility management as a
high-priority use case. Bonati \textit{et~al.}~\cite{bonati_colosseum}
demonstrate an end-to-end RL xApp on the Colosseum testbed. SON-GNN-DQN
is designed to be xApp-compatible: the GNN-DQN inference fits within the
$10$--$1000$~ms Near-RT loop, and the SON translation layer outputs an
action format consistent with E2SM-RC control messages.

\section{Drive-Test-Based Evaluation in Developing Markets}

The literature on drive-test validation in South Asian contexts is
sparse but growing. Khanal \textit{et~al.}~\cite{khanal_nepal}
characterize LTE coverage along the Prithvi Highway, and Bhattarai
\textit{et~al.}~\cite{bhattarai_kathmandu} study handover failures in
Kathmandu Valley. Both works call for AI-driven optimization and
provide empirical evidence that handover anomalies in dense Nepalese
clusters are non-trivial. The validation phase of the present project is
aligned with this body of evidence.

\section{Research Gap}

Three converging gaps motivate this project. First, RL-based handover
policies have been demonstrated in simulation but rarely validated
against operator-grade data from a developing-country LTE deployment.
Second, although GNNs are recognized as natural representations for
cellular networks, their combination with DQN under a SON-translation
envelope has received limited attention. Third, the bridge between a
research model and operator-ready SON parameters (bounded CIO and TTT
updates) is rarely made explicit. This project addresses these gaps by
proposing a topology-generalized GNN-DQN agent wrapped in a safety-
bounded SON translation layer; the planned validation phase will close
the data gap through field measurements and, subject to NTC's
cooperation, operator-side parameters.

%==============================================================================
\chapter{RELATED THEORY}
%==============================================================================

\section{LTE Network Fundamentals}

LTE is based on the Evolved Universal Terrestrial Radio Access Network
(E-UTRAN) and the Evolved Packet Core (EPC). The UE communicates with an
eNodeB (eNB), which manages radio resource control, scheduling and
handovers. The EPC includes the Mobility Management Entity (MME), the
Serving Gateway (SGW), the Packet Gateway (PGW) and the PCRF. Within a
cell, LTE allocates resources on a Physical Resource Block (PRB) basis,
where each PRB consists of $12$ subcarriers across a $0.5$~ms slot. PRB
utilization is therefore a direct measure of cell load. The downlink
throughput available to a UE is approximately proportional to the
fraction of unused PRBs and to the Shannon capacity dictated by the
achieved Signal-to-Interference-plus-Noise Ratio (SINR).

\begin{figure}[H]
    \centering
    \includegraphics[width=0.85\textwidth]{figures/lte_arch.png}
    \caption{LTE network architecture showing E-UTRAN and EPC components.}
    \label{fig:lte_arch}
\end{figure}

\section{Reference Signal Measurements}

The two principal UE-side measurements used for mobility decisions are:
\begin{itemize}[leftmargin=*]
    \item \textbf{RSRP} (Reference Signal Received Power): the linear
    average of the power contributions of resource elements that carry
    cell-specific reference signals. Typical values range from $-44$~dBm
    (very strong) to $-140$~dBm (out of coverage).
    \item \textbf{RSRQ} (Reference Signal Received Quality), defined as
    $\mathrm{RSRQ} = N \cdot \mathrm{RSRP}/\mathrm{RSSI}$, captures both
    signal strength and interference / load. Lightly loaded cells exhibit
    RSRQ in the range $-3$ to $-6$~dB, while heavily loaded cells degrade
    to $-15$~dB or worse. RSRQ is therefore a practical UE-side load
    proxy when PRB utilization is unavailable.
\end{itemize}
The SINR completes the picture by capturing the joint effect of
interference from co-channel cells and noise.

\section{LTE Handover Procedure and Event A3}

Handover in LTE is a network-controlled, UE-assisted procedure. The
serving eNB configures the UE with measurement objects and reporting
events. When the Event~A3 condition,
\begin{equation}
    \mathrm{RSRP}_{\text{neighbor}} + \mathrm{CIO}_{\text{neighbor}}
    \;\geq\;
    \mathrm{RSRP}_{\text{serving}} + \mathrm{CIO}_{\text{serving}}
    + A3_{\text{offset}},
\end{equation}
holds for the entire Time-to-Trigger (TTT) duration, the UE sends a
Measurement Report. The source eNB decides on handover, initiates an X2
Handover Request to the target eNB, receives an acknowledgement, and
issues an RRC Connection Reconfiguration message that instructs the UE
to perform random access on the target cell. Path switching with the
MME/SGW finalises the procedure.

\begin{figure}[H]
    \centering
    \includegraphics[width=0.9\textwidth]{figures/a3_event.png}
    \caption{Event~A3 trigger and the associated handover signalling
    sequence in LTE.}
    \label{fig:a3_event}
\end{figure}

\section{Cell Individual Offset and SON Translation}

The Cell Individual Offset (CIO) is a per-neighbour parameter
($\pm 24$~dB, $0.5$~dB resolution) that biases the A3 inequality. A
positive CIO toward a neighbour cell makes that cell appear stronger,
encouraging earlier handover; a negative CIO discourages it. SON
Mobility Load Balancing adjusts CIOs on a slow time scale to redistribute
UEs across cells. In SON-GNN-DQN, the per-UE target preference learned
by the agent is translated into a CIO update through
\begin{equation}
    \Delta\mathrm{CIO}_{i \to j} \;=\;
    \mathrm{clip}\!\Big(
    \alpha\cdot\big(Q(s, j) - Q(s, i)\big),
    -\Delta_{\max}, +\Delta_{\max}
    \Big),
    \label{eq:cio_translation}
\end{equation}
where $Q(s,\cdot)$ are the Q-values from the DQN, $\alpha$ is a scaling
factor, and $\Delta_{\max} = 3$~dB per minute enforces a SON safety
bound. The cumulative CIO is constrained to remain within $\pm 9$~dB to
keep the network within a conservative operating envelope.

\section{Reinforcement Learning and Deep Q-Networks}

Reinforcement Learning (RL) formalises the interaction between an agent
and an environment as a Markov Decision Process
$(\mathcal{S}, \mathcal{A}, \mathcal{P}, \mathcal{R}, \gamma)$. The
agent selects actions $a_t \in \mathcal{A}$ in state $s_t$ to maximize
the expected discounted return
$G_t = \sum_{k=0}^{\infty}\gamma^{k} r_{t+k+1}$.
The action-value function $Q^{\pi}(s,a)$ satisfies the Bellman equation
\begin{equation}
    Q^{\pi}(s,a) = \mathbb{E}\!\left[ r_{t+1} + \gamma\,
    Q^{\pi}(s_{t+1}, a_{t+1}) \mid s_t=s, a_t=a \right].
\end{equation}
Deep Q-Networks parameterise $Q(s,a;\theta)$ with a neural network and
stabilise training through experience replay, target networks and an
$\epsilon$-greedy exploration policy \cite{mnih_dqn_2015}. Double DQN
\cite{vanhasselt_doubledqn} reduces over-estimation bias and is used in
this project together with a dueling head \cite{wang_dueling_2016} and
a prioritized replay buffer \cite{schaul_per_2016}.

\section{Graph Neural Networks}

A graph $G = (\mathcal{V}, \mathcal{E})$ with node features
$\mathbf{h}^{(0)}_i$ is processed through $L$ message-passing layers in
which each node aggregates information from its one-hop neighbourhood:
\begin{equation}
    \mathbf{h}^{(\ell+1)}_i = \sigma\!\Big( W_1 \mathbf{h}^{(\ell)}_i +
    W_2 \!\! \sum_{j \in \mathcal{N}(i)} \!\! \alpha_{ij}\,
    \mathbf{h}^{(\ell)}_j \Big),
    \label{eq:gnn_msg}
\end{equation}
where $\alpha_{ij}$ is an attention coefficient as in Graph Attention
Networks (GAT) \cite{velickovic_gat}. In SON-GNN-DQN the GNN is a two-
layer GAT with $4$ attention heads and $64$-dimensional hidden states.
Because the message-passing rule is permutation-equivariant, the GNN can
process graphs of different sizes without retraining, which supports
the topology-generalization requirement.

\begin{figure}[H]
    \centering
    \includegraphics[width=0.78\textwidth]{figures/gnn_dqn_arch.png}
    \caption{High-level architecture of SON-GNN-DQN. The GNN consumes the
    dynamic cell graph and emits per-cell embeddings; the DQN consumes
    the UE state concatenated with the embeddings and emits Q-values
    that drive the SON translation layer.}
    \label{fig:gnn_dqn_arch}
\end{figure}

%==============================================================================
\chapter{SYSTEM ARCHITECTURE}
%==============================================================================

\section{Overall System Block Diagram}

SON-GNN-DQN is organised as a closed-loop pipeline comprising four
functional layers: a measurement and data ingestion layer, a graph
construction layer, a learning core (GNN encoder and DQN agent), and a
SON translation and execution layer. Each layer is decoupled through
clean data contracts so that the synthetic simulator and (in the
validation phase) real drive-test or operator data can be processed by
the same downstream stack.

\begin{figure}[H]
    \centering
    \includegraphics[width=0.95\textwidth]{figures/system_block.png}
    \caption{Overall system block diagram of SON-GNN-DQN, showing the
    measurement and ingestion layer, the graph construction layer, the
    GNN-DQN learning core, and the safety-bounded SON translation
    layer.}
    \label{fig:system_block}
\end{figure}

\section{Data Ingestion Layer}

In this interim phase the only active source is the synthetic LTE
simulator, which produces RSRP, RSRQ, SINR, throughput, UE position,
velocity, PRB load and handover events at $1$~Hz over a configurable
topology. To prepare for the validation phase, the data ingestion layer
already supports a second source: cleaned drive-test logs collected on
an Android phone (Poco~X3) with the Network Cell Info Lite application,
locked to LTE only. Both sources are mapped to a unified CSV schema
with the columns \texttt{timestamp, lat, lon, speed, serving\_PCI,
RSRP, RSRQ, SINR, neighbor\_PCI\_k, neighbor\_RSRP\_k} for
$k = 1, \ldots, 6$, so that the same preprocessing pipeline can operate
on either stream.

\section{Graph Construction Layer}

For every timestep $t$ a dynamic graph
$G_t = (\mathcal{V}_t, \mathcal{E}_t)$ is constructed.
\begin{itemize}[leftmargin=*]
    \item \textbf{Nodes.} Each visible cell (unique Physical Cell
    Identifier in the serving and neighbour set) becomes a node. The UE
    is represented as a special node connected to all cell nodes.
    \item \textbf{Node features.} For each cell node the
    \texttt{ue\_only} profile uses
    $\{\mathrm{RSRP}, \mathrm{RSRQ},$ RSRQ-derived load proxy,
    $\mathrm{RSRP}$ trend, $\mathrm{RSRQ}$ trend, serving-cell
    indicator, previous serving indicator, signal usability, UE speed,
    time since last handover$\}$.
    \item \textbf{Edges.} Two cells are connected if they appear together
    in the UE's neighbour list within any $5$-second window. Edge
    weights are the absolute RSRP difference between the connected cells.
\end{itemize}
This yields a $10$-dimensional feature vector per cell node and a $5$-
dimensional feature vector for the UE node. An extended
\texttt{oran\_e2} profile reserves slots for $\mathrm{PRB}_{\text{util}}$,
$\mathrm{PRB}_{\text{avail}}$ and connected-UE count; these are
populated only when operator-side data become available.

\section{GNN-DQN Learning Core}

The learning core consists of a two-layer Graph Attention Network
followed by a Double + Dueling DQN head. The GAT produces $64$-
dimensional embeddings per cell node; the embeddings are concatenated
with the UE node embedding and passed to a dueling head that decomposes
$Q(s,a)$ into a value stream $V(s)$ and an advantage stream $A(s,a)$.
Action selection uses an $\epsilon$-greedy policy with linear decay from
$\epsilon = 1.0$ to $\epsilon = 0.05$ over the first $200{,}000$
environment steps. The replay buffer is a prioritized buffer of size
$200{,}000$ with $\alpha = 0.6$ and $\beta$ annealed from $0.4$ to $1.0$.

\section{SON Translation and Execution Layer}

The SON translation layer is the bridge between the learning core and a
future deployment. It receives the soft preference scores from the
GNN-DQN, applies Equation~\eqref{eq:cio_translation} to compute bounded
CIO updates, and writes the updates either to the simulator (for
evaluation) or to a JSON record consistent with the E2SM-RC action
format. Three SON modes are supported: \texttt{shadow} (log-only),
\texttt{advise} (applies CIO updates but disables TTT changes) and
\texttt{full} (applies both CIO and TTT updates within the safety
bounds). All experiments in this interim phase use \texttt{advise} mode.

%==============================================================================
\chapter{METHODOLOGY}
%==============================================================================

\section{Overall Workflow}

The methodology proceeds in five stages: (i) environment design and
synthetic data generation; (ii) feature engineering and graph
construction; (iii) implementation of the GNN-DQN agent, baselines and
SON translation layer; (iv) staged training using a curriculum that
progressively adds urban, highway and hilly scenarios; and (v) multi-
method evaluation under multiple seeds. Stage (vi), drive-test
validation in Pokhara and operator-side parameter validation with NTC,
is the next phase of the project and is described separately in
Chapter~9.

\section{Synthetic Simulation Environment}

The simulator is implemented in Python using NumPy and PyTorch tensors
for vectorised computation. It supports configurable topologies of
$4$ to $19$ base stations on a $5 \times 5$~km$^2$ area. Path loss
follows the log-distance model
$\mathrm{PL}(d) = \mathrm{PL}_0 + 10\,n\,\log_{10}(d) + X_\sigma$, with
$n \in \{2.6, 3.5, 4.2\}$ for urban, highway and hilly profiles and a
log-normal shadowing standard deviation
$\sigma \in \{4, 6, 9\}$~dB. SINR is computed with explicit interference
aggregation over all visible cells on the same frequency. Throughput is
derived through the Shannon expression scaled by the fraction of unused
PRBs in the serving cell. PRB utilization evolves through a Poisson
arrival model. UE mobility is modelled using a Gauss--Markov process
for direction and a uniform distribution for speed within a scenario-
specific range. Three scenarios are defined for training and
evaluation, summarised in Table~\ref{tab:scenarios}.

\begin{table}[H]
    \centering
    \caption{Synthetic training scenarios.}
    \label{tab:scenarios}
    \small
    \begin{tabularx}{\textwidth}{lcccc}
        \toprule
        \textbf{Scenario} & \textbf{Cells} & \textbf{Inter-site} &
        \textbf{UE Speed} & \textbf{$n$ / $\sigma$ (dB)} \\
        \midrule
        Urban Dense  & 12 & 350~m  & 25--40~km/h & 3.5 / 6.0\\
        Highway      &  6 & 1500~m & 50--80~km/h & 2.6 / 4.0\\
        Hilly Shadow &  5 & 1800~m & 20--35~km/h & 4.2 / 9.0\\
        \bottomrule
    \end{tabularx}
\end{table}

\section{Feature Engineering and Normalisation}

All continuous features are standardised using the per-scenario mean
and standard deviation computed on the training split, with values
clipped to $[-3, +3]$ before being passed to the GNN. The RSRQ-derived
load proxy is computed as
\begin{equation}
    \tilde{\rho}_i = \mathrm{clip}\!\left(
    \frac{-3 - \mathrm{RSRQ}_i}{16.5}, 0.05, 0.95\right),
\end{equation}
which maps $\mathrm{RSRQ} = -3$~dB to $\tilde{\rho} = 0.1$ and
$\mathrm{RSRQ} = -19.5$~dB to $\tilde{\rho} = 0.95$. The signal-
usability feature is computed as
$\sigma_u(\mathrm{RSRP}) = \tanh\!\big((\mathrm{RSRP}+95)/12\big)$.

\section{Reward Function Design}

The reward at step $t$ is a weighted sum of throughput gain, ping-pong
penalty, RLF penalty, load-imbalance penalty and handover cost:
\begin{equation}
\begin{aligned}
r_t \;=\;& w_{1}\,\Delta\widehat{\text{Tput}}_t \;
        - w_{2}\,\mathbb{1}_{\text{ping-pong}} \;
        - w_{3}\,\mathbb{1}_{\text{RLF}}\\
        &- w_{4}\,J(\mathbf{L}_t)^{-1} \;
        - w_{5}\,\mathbb{1}_{\text{handover}},
\end{aligned}
\end{equation}
where $J(\mathbf{L}_t)$ is Jain's fairness index over PRB-load vectors
and $\mathbb{1}_{\text{ping-pong}}$ is set when the same UE returns to
the previous serving cell within $5$~s. After three rounds of reward
shaping the current weights are $w_1=1.0$, $w_2=2.5$, $w_3=10.0$,
$w_4=0.8$, $w_5=0.2$. The handover-cost penalty is deliberately small
to avoid an overly conservative agent.

\section{Training Procedure}

Training is organised as a curriculum that gradually introduces more
difficult scenarios, summarised in Table~\ref{tab:curriculum}.

\begin{table}[H]
    \centering
    \caption{Curriculum training schedule. The smoke and diagnostic
    stages have been completed; the multi-scenario long run is in
    progress at the time of this interim report.}
    \label{tab:curriculum}
    \small
    \begin{tabularx}{\textwidth}{llcl}
        \toprule
        \textbf{Stage} & \textbf{Scenario(s)} & \textbf{Episodes} &
        \textbf{Status}\\
        \midrule
        Smoke (UE-only)         & Urban only        &  Short
        & Completed (sanity check passed)\\
        Diagnostic (UE-only)    & Urban + Highway   & Medium
        & Completed (return curve rises)\\
        Multi-scenario          & Urban + Highway + Hilly & Long
        & In progress (Google Colab Pro GPU)\\
        Drive-test fine-tune    & Real Pokhara routes & --
        & Planned (next phase)\\
        \bottomrule
    \end{tabularx}
\end{table}

Training is executed on a personal laptop (Apple MacBook Pro M1 Pro,
16~GB unified memory, 512~GB storage) for development and on Google
Colab Pro (NVIDIA T4 / A100 depending on availability) for the longer
multi-scenario run. Mixed-precision (FP16) arithmetic is used for the
GAT forward pass on Colab to reduce memory pressure; the DQN head
remains in FP32 for numerical stability of the Bellman update. Resume
checkpoints are kept lightweight (about $28$~MB); a full replay
snapshot is enabled only at the end of each stage because of its size
(approximately $2.1$~GB).

\section{Pseudo-Code}

\begin{algorithm}[H]
\caption{SON-GNN-DQN training procedure}
\label{alg:train}
\begin{algorithmic}[1]
\State Initialize GAT $\phi$, dueling DQN $Q_\theta$, target
$Q_{\theta^-}$, prioritized replay $B$
\For{stage $\in$ \{Smoke, Diagnostic, Multi, Drive-test\}}
  \For{episode $= 1, \ldots, N_{\text{stage}}$}
    \State $s_0 \gets$ env.reset(stage)
    \For{$t = 0, 1, \ldots, T-1$}
      \State $G_t \gets$ ConstructGraph($s_t$)
      \State $\mathbf{e}_t \gets \phi(G_t)$
      \State $x_t \gets [\mathbf{e}_t,\, \text{UE-features}]$
      \State $a_t \gets \arg\max_a Q_\theta(x_t, a)$ with probability
              $1-\epsilon$, else random
      \State $(s_{t+1}, r_t, \text{done}) \gets$ env.step($a_t$)
      \State Apply SON translation \eqref{eq:cio_translation}
      (advise mode)
      \State Push $(s_t, a_t, r_t, s_{t+1})$ into $B$
      \State Sample minibatch from $B$ with priorities
      \State Compute Double-DQN target
      $y = r + \gamma\, Q_{\theta^-}\!\big(s', \arg\max_{a'}
      Q_\theta(s', a')\big)$
      \State Update $\theta$ via gradient descent on
      $\big(Q_\theta(s,a) - y\big)^2$
      \State Every $C$ steps: $\theta^- \gets \theta$
    \EndFor
  \EndFor
\EndFor
\State \Return $\theta, \phi$
\end{algorithmic}
\end{algorithm}

\section{Evaluation Methods (Planned)}

In the final phase the agent will be evaluated against six baselines
plus its own SON-translated variant: \texttt{no\_handover},
\texttt{random\_valid}, \texttt{strongest\_rsrp}, \texttt{a3\_ttt}
(default $A3 = 3$~dB, $\mathrm{TTT} = 160$~ms),
\texttt{load\_aware} (a static-threshold MLB variant), \texttt{gnn\_dqn}
(raw learned policy) and \texttt{son\_gnn\_dqn} (the SON-translated,
deployment-friendly variant). Each method will be run over multiple
seeds and reported with $95\%$ confidence intervals. The baseline
implementations are already in the repository and have been smoke-
tested; the final-table numbers will be produced once the multi-
scenario training run and the drive-test campaign are complete.

%==============================================================================
\chapter{IMPLEMENTATION}
%==============================================================================

\section{Software Stack}

The project is implemented entirely in Python~3.11 with PyTorch~2.2 and
PyTorch~Geometric~2.5 for graph operations. Supporting libraries include
NumPy~1.26, pandas~2.2, SciPy~1.13, scikit-learn~1.4, matplotlib~3.8 and
seaborn~0.13. The simulator follows an OpenAI Gym-compatible interface
with \texttt{reset()} and \texttt{step()} signatures. Reproducibility is
ensured by fixed seeds for Python, NumPy and PyTorch (CPU and CUDA),
and by recording the exact configuration file used for each run inside
the checkpoint metadata.

\section{Hardware Used}

\begin{itemize}[leftmargin=*]
    \item \textbf{Development laptop}: Apple MacBook Pro (M1 Pro, 2020),
    16~GB unified memory, 512~GB SSD. Used for code development,
    smoke runs, diagnostic runs, and CPU-side data processing.
    \item \textbf{Cloud GPU}: Google Colab Pro (NVIDIA T4 / A100,
    depending on availability), used for the multi-scenario long
    training run.
    \item \textbf{Drive-test handset (pilot only)}: Poco~X3 (Android),
    LTE-only locked, Network Cell Info Lite at $1$~Hz. The pilot
    sessions verified end-to-end logging; representative drive-test
    runs are part of the next phase.
\end{itemize}

\section{Repository Structure}

The repository follows the layout shown in Listing~\ref{lst:repo}. Top-
level scripts (\texttt{scripts/train.py}, \texttt{scripts/evaluate.py})
are kept thin and delegate to library modules.

\begin{lstlisting}[language=bash,caption={Repository layout (top-level
folders).},label={lst:repo}]
gnn-dqn-handover-loadbalancing/
+-- configs/experiments/*.json    # smoke, diagnostic, multiscenario, ...
+-- scripts/                       # train.py, evaluate.py, drive_test.py
+-- src/
|   +-- env/                       # LTE simulator, mobility, path loss
|   +-- graph/                     # graph construction, normalisation
|   +-- agent/                     # GAT encoder, dueling DQN, replay
|   +-- son/                       # translation layer, CIO/TTT bounds
|   +-- baselines/                 # a3_ttt, load_aware, random_valid, ...
|   +-- eval/                      # metrics, plotting, multi-seed runner
+-- results/runs/                   # multiscenario_ue/, smoke_ue/, ...
+-- tests/                          # pytest suite, smoke checks
\end{lstlisting}

\section{Reproducing the Current Pipeline}

The current pipeline can be reproduced with the following commands:

\begin{lstlisting}[language=bash,caption={Commands used during this
interim phase.}]
PYTHONPATH=src python3 -m pytest -q
python3 scripts/train.py --config configs/experiments/smoke_ue.json
python3 scripts/train.py --config configs/experiments/diagnostic_ue.json
python3 scripts/train.py --config configs/experiments/multiscenario_ue.json
# Multi-scenario run is long; intermediate resume checkpoints are
# saved under results/runs/multiscenario_ue/checkpoints/resume/.
\end{lstlisting}

%==============================================================================
\chapter{PRELIMINARY RESULTS}
%==============================================================================

This chapter reports only verified outcomes of the work completed so far.
Numerical results from the long multi-scenario run and from the drive-
test campaign are intentionally deferred to the final submission. The
purpose of this chapter is to demonstrate that the pipeline is working
end-to-end and that the agent shows the qualitative learning behaviour
expected from the design.

\section{What Has Been Verified}

\begin{itemize}[leftmargin=*]
    \item The synthetic simulator produces RSRP, RSRQ, SINR and
    throughput trajectories that are internally consistent (units,
    dynamic range, sign conventions, log-distance path-loss curve).
    \item The graph construction module produces valid PyTorch Geometric
    \texttt{Data} objects on every step under all three scenario types.
    Batching with \texttt{Batch.from\_data\_list} works correctly for
    variable-size graphs.
    \item The GAT encoder, the Double + Dueling DQN head and the
    prioritized replay buffer integrate cleanly: gradient flow has been
    verified by gradient-norm checks during the smoke and diagnostic
    runs.
    \item All six baselines (\texttt{no\_handover}, \texttt{random\_valid},
    \texttt{strongest\_rsrp}, \texttt{a3\_ttt}, \texttt{load\_aware},
    \texttt{gnn\_dqn}, \texttt{son\_gnn\_dqn}) execute against the same
    simulator interface.
    \item The SON translation layer enforces the bounded CIO update
    rule (Equation~\eqref{eq:cio_translation}); a unit test in
    \texttt{tests/test\_son\_bounds.py} confirms that the cumulative
    CIO is clipped to $\pm 9$~dB across long episodes.
    \item The unit-test suite (\texttt{pytest -q}) passes on the current
    revision.
\end{itemize}

\section{Qualitative Behaviour Observed in Smoke / Diagnostic Runs}

Figure~\ref{fig:training_curve_qual} shows a qualitative sketch of the
return trajectories observed in the smoke and diagnostic runs. In both
runs the agent crosses the \texttt{random\_valid} baseline within a few
tens of episodes and continues to improve as the curriculum advances.
Specific numerical values are not reported here because the diagnostic
run is still preliminary; the final report will replace the sketch
with the actual logged curves from the multi-scenario run together
with confidence intervals.

\begin{figure}[H]
\centering
\begin{tikzpicture}
\begin{axis}[
    width=0.85\textwidth, height=6.5cm,
    xlabel={Episode (illustrative)},
    ylabel={Return (illustrative)},
    xmin=0, xmax=100, ymin=-30, ymax=80,
    grid=both, grid style={line width=.1pt, draw=gray!30},
    legend pos=south east, legend cell align=left,
    every axis plot/.append style={thick}
]
\addplot[blue] coordinates {(0,-25) (5,-10) (10,-2) (20,12)
    (30,24) (40,34) (50,42) (60,48) (70,54) (80,58) (90,60) (100,62)};
\addplot[red,dashed] coordinates {(0,-15) (100,-15)};
\legend{SON-GNN-DQN (smoke run, qualitative), Random-valid baseline}
\end{axis}
\end{tikzpicture}
\caption{Qualitative trend of the return observed in the smoke and
diagnostic runs of the SON-GNN-DQN agent. Numerical values are not
intended to be quoted; the final report will replace this figure with
fully logged curves over multiple seeds.}
\label{fig:training_curve_qual}
\end{figure}

\section{What is Pending}

\begin{itemize}[leftmargin=*]
    \item Completion of the multi-scenario training run and multi-seed
    evaluation.
    \item Quantitative comparison against all baselines with $95\%$
    confidence intervals.
    \item Drive-test campaign in Pokhara (Lakeside--Mahendrapul,
    Prithvi Highway and Bindabasini--Sarangkot routes).
    \item Validation against NTC operator data (subject to NTC's
    cooperation under the request letter in Appendix~D).
\end{itemize}

%==============================================================================
\chapter{CHALLENGES ENCOUNTERED SO FAR}
%==============================================================================

\textbf{Reward shaping.} Early reward formulations either rewarded too
many handovers (the agent oscillated) or penalised them too aggressively
(the agent froze on a single cell). Three rounds of reward shaping
were required to reach the current weights.

\textbf{Dynamic graph batching.} Each timestep produces a graph of
variable size, which complicates batching across episodes. The project
relies on PyTorch Geometric's \texttt{Batch.from\_data\_list} to
disjoint-union the graphs and uses custom collate functions to keep
replay-buffer samples graph-aware.

\textbf{Hardware constraints.} The development laptop (M1 Pro, 16~GB)
cannot host a full multi-scenario training run alongside a development
workflow. The pragmatic compromise is to develop and unit-test locally
and to run the longer experiments on Google Colab Pro, with
intermediate checkpoints synchronised to a private repository.

\textbf{Drive-test pilot.} The pilot logging session on the Poco~X3
revealed that battery-saver mode and Wi-Fi auto-switch silently drop
log rows. A pre-trip checklist has been written (see Appendix~B) to
prevent these issues during the upcoming representative campaign.

%==============================================================================
\chapter{REMAINING WORK AND TIMELINE}
%==============================================================================

\section{Remaining Work}

\begin{enumerate}[leftmargin=*]
    \item \textbf{Complete the multi-scenario long training run} and
    persist the final checkpoint and replay snapshot.
    \item \textbf{Run multi-seed evaluation} of all baselines and the
    SON-GNN-DQN policy on a held-out synthetic test split.
    \item \textbf{Conduct the structured interview with NTC engineers}
    in Pokhara using the question list in Appendix~C.
    \item \textbf{Request anonymised drive-test, OSS and configuration
    data} from NTC under the formal letter in Appendix~D and integrate
    any approved data into the unified pipeline.
    \item \textbf{Carry out a representative drive-test campaign in
    Pokhara} on the three planned routes (Lakeside--Mahendrapul,
    Prithvi Highway, Bindabasini--Sarangkot).
    \item \textbf{Fine-tune the agent} on the integrated real-world
    data and re-run the evaluation pipeline.
    \item \textbf{Produce final figures and tables} with the actual
    numbers and update the report.
\end{enumerate}

\section{Indicative Timeline}

\begin{table}[H]
    \centering
    \caption{Indicative timeline for the remaining work.}
    \small
    \begin{tabularx}{\textwidth}{lX}
        \toprule
        \textbf{Week} & \textbf{Activity}\\
        \midrule
        Week 1--2 & Complete multi-scenario long training run; run
        multi-seed evaluation.\\
        Week 3   & Conduct NTC engineer interview; submit formal data
        request.\\
        Week 4--5 & Drive-test campaign in Pokhara; data cleaning and
        ingestion.\\
        Week 6   & Fine-tuning and re-evaluation.\\
        Week 7   & Update report, finalise figures, prepare defence.\\
        \bottomrule
    \end{tabularx}
\end{table}

%==============================================================================
\chapter{CONCLUSION (INTERIM)}
%==============================================================================

This interim progress report describes the design and implementation of
SON-GNN-DQN, a topology-generalized Graph Neural Network and Deep
Q-Learning framework wrapped in a safety-bounded SON translation layer
for adaptive handover and load optimization in LTE networks. The
synthetic LTE simulator, the unified CSV schema, the graph-construction
module, the GAT encoder, the Double + Dueling DQN agent with prioritized
replay, the six baseline policies, the SON translation layer with the
bounded CIO update rule, and the unit-test suite are in place and have
been verified through smoke and diagnostic runs. The multi-scenario
long training run is in progress on Google Colab Pro at the time of
writing. The drive-test campaign in Pokhara and the validation against
Nepal Telecom operator data are the next steps and are documented in
the remaining-work chapter and the data-request letter in Appendix~D.

%==============================================================================
\chapter*{REFERENCES}
\addcontentsline{toc}{chapter}{REFERENCES}
%==============================================================================

\begin{thebibliography}{99}

\bibitem{ts36331} 3GPP, ``Evolved Universal Terrestrial Radio Access
(E-UTRA); Radio Resource Control (RRC); Protocol specification,''
TS~36.331, Release~17, 2023.

\bibitem{3gpp_36839} 3GPP, ``Mobility enhancements in heterogeneous
networks,'' TR~36.839, Release~11, 2012.

\bibitem{ts36902} 3GPP, ``Self-configuring and self-optimizing network
(SON) use cases and solutions,'' TR~36.902, Release~9, 2011.

\bibitem{oran_alliance_wg2} O-RAN Alliance, ``O-RAN Working Group~2:
Non-RT RIC and A1 Interface Architecture,'' Technical Specification,
O-RAN.WG2.AAD, v.06.00, 2023.

\bibitem{polese_oran_survey} M.~Polese, L.~Bonati, S.~D'Oro, S.~Basagni,
and T.~Melodia, ``Understanding O-RAN: Architecture, interfaces,
algorithms, security, and research challenges,'' \textit{IEEE
Communications Surveys \& Tutorials}, vol.~25, no.~2, pp.~1376--1411,
2023.

\bibitem{bonati_colosseum} L.~Bonati \textit{et~al.}, ``Intelligence
and learning in O-RAN for data-driven NextG cellular networks,''
\textit{IEEE Communications Magazine}, vol.~59, no.~10, pp.~21--27,
2021.

\bibitem{mnih_dqn_2015} V.~Mnih \textit{et~al.}, ``Human-level control
through deep reinforcement learning,'' \textit{Nature}, vol.~518,
no.~7540, pp.~529--533, 2015.

\bibitem{vanhasselt_doubledqn} H.~van~Hasselt, A.~Guez, and D.~Silver,
``Deep reinforcement learning with double Q-learning,'' in
\textit{Proc.\ AAAI}, 2016, pp.~2094--2100.

\bibitem{wang_dueling_2016} Z.~Wang \textit{et~al.}, ``Dueling network
architectures for deep reinforcement learning,'' in \textit{Proc.\
ICML}, 2016.

\bibitem{schaul_per_2016} T.~Schaul, J.~Quan, I.~Antonoglou, and
D.~Silver, ``Prioritized experience replay,'' in \textit{Proc.\
ICLR}, 2016.

\bibitem{kipf_gcn} T.~N.~Kipf and M.~Welling, ``Semi-supervised
classification with graph convolutional networks,'' in
\textit{Proc.\ ICLR}, 2017.

\bibitem{velickovic_gat} P.~Veli\v{c}kovi\'c \textit{et~al.}, ``Graph
attention networks,'' in \textit{Proc.\ ICLR}, 2018.

\bibitem{hamilton_graphsage} W.~Hamilton, R.~Ying, and J.~Leskovec,
``Inductive representation learning on large graphs,'' in
\textit{Proc.\ NeurIPS}, 2017.

\bibitem{eisen_random_graphs_2020} M.~Eisen and A.~Ribeiro, ``Optimal
wireless resource allocation with random edge graph neural networks,''
\textit{IEEE Transactions on Signal Processing}, vol.~68,
pp.~2977--2991, 2020.

\bibitem{shen_gnn_resource_2021} Y.~Shen, Y.~Shi, J.~Zhang, and
K.~B.~Letaief, ``Graph neural networks for scalable radio resource
management: Architecture design and theoretical analysis,''
\textit{IEEE Journal on Selected Areas in Communications}, vol.~39,
no.~1, pp.~101--115, 2021.

\bibitem{tang_handover_drl} F.~Tang, B.~Mao, Y.~Kawamoto, and N.~Kato,
``Survey on machine learning for intelligent end-to-end communication
toward 6G: From network access, routing to traffic control and streaming
adaption,'' \textit{IEEE Communications Surveys \& Tutorials},
vol.~23, no.~3, pp.~1578--1598, 2021.

\bibitem{wang_handover_drl_2022} W.~Wang, J.~Yao, L.~Zhang, and Z.~Han,
``DRL-based intelligent handover for high-mobility scenarios in dense
cellular networks,'' \textit{IEEE Transactions on Vehicular
Technology}, vol.~71, no.~5, pp.~5302--5315, 2022.

\bibitem{zhao_gnn_handover} R.~Zhao, J.~Lu, and W.~Zhang, ``A graph-
neural-network-aided deep reinforcement learning framework for handover
optimisation in heterogeneous networks,'' \textit{IEEE Transactions on
Wireless Communications}, vol.~23, no.~4, pp.~2701--2715, 2024.

\bibitem{wang_gnn_rl_2022} N.~Wang, J.~Tang, and K.~K.~Wong, ``Graph-
based deep reinforcement learning for joint user association and
resource allocation in dense networks,'' \textit{IEEE Transactions on
Communications}, vol.~70, no.~9, pp.~6184--6198, 2022.

\bibitem{he_vanet_gnn} Y.~He \textit{et~al.}, ``Graph neural network-
based routing in vehicular ad hoc networks,'' \textit{IEEE Internet of
Things Journal}, vol.~9, no.~13, pp.~10538--10551, 2022.

\bibitem{castro_handover_2014} L.~Castro-Hernandez and R.~Paranthaman,
``Optimization of handover parameters in heterogeneous networks,''
in \textit{Proc.\ IEEE WCNC}, 2014, pp.~2820--2825.

\bibitem{lopezperez_mobility_2012} D.~Lopez-Perez, I.~Guvenc, and
X.~Chu, ``Mobility management challenges in 3GPP heterogeneous
networks,'' \textit{IEEE Communications Magazine}, vol.~50, no.~12,
pp.~70--78, 2012.

\bibitem{zhang_mlb_2021} J.~Zhang \textit{et~al.}, ``A survey of
mobility load balancing for self-organizing cellular networks,''
\textit{IEEE Communications Surveys \& Tutorials}, vol.~23, no.~4,
pp.~2347--2376, 2021.

\bibitem{hasan_son_2020} M.~Hasan and A.~Mohamed, ``Fuzzy logic-based
joint MRO and MLB in LTE networks,'' in \textit{Proc.\ IEEE GLOBECOM},
2020.

\bibitem{khanal_nepal} S.~Khanal, B.~Adhikari, and R.~Pandey, ``LTE
coverage characterisation along the Prithvi Highway,'' \textit{Journal
of the Institute of Engineering (JIE)}, vol.~17, no.~1, pp.~88--95,
2023.

\bibitem{bhattarai_kathmandu} S.~Bhattarai and S.~Khanal, ``Empirical
study of handover failures in Kathmandu Valley LTE deployments,''
\textit{IOE Graduate Conference Proceedings}, 2024.

\bibitem{fey_pyg_2019} M.~Fey and J.~E.~Lenssen, ``Fast graph
representation learning with PyTorch Geometric,'' in \textit{Proc.\
ICLR Workshop on Representation Learning on Graphs and Manifolds},
2019.

\bibitem{paszke_pytorch_2019} A.~Paszke \textit{et~al.}, ``PyTorch: An
imperative style, high-performance deep learning library,'' in
\textit{Proc.\ NeurIPS}, 2019.

\end{thebibliography}

%==============================================================================
\appendix
%==============================================================================

\chapter{HYPERPARAMETERS USED SO FAR}

\begin{table}[H]
    \centering
    \caption{Current hyperparameters of the GNN-DQN agent and SON
    translation layer.}
    \footnotesize
    \begin{tabularx}{0.85\textwidth}{Xl}
        \toprule
        \textbf{Hyperparameter} & \textbf{Value}\\
        \midrule
        GAT layers                  & 2\\
        GAT hidden size             & 64\\
        GAT attention heads         & 4\\
        GAT dropout                 & 0.1\\
        DQN hidden layers           & [128, 128]\\
        DQN type                    & Double + Dueling\\
        Learning rate (Adam)        & $3 \times 10^{-4}$\\
        Discount $\gamma$           & 0.97\\
        Soft target update $\tau$   & 0.005\\
        Batch size                  & 64\\
        Replay buffer               & 200{,}000 (prioritized)\\
        $\epsilon$ schedule         & $1.0 \to 0.05$ linear over
        $2{\times}10^{5}$ steps\\
        SON $\Delta_{\max}$/update  & 3~dB\\
        SON cumulative CIO bound    & $\pm 9$~dB\\
        \bottomrule
    \end{tabularx}
\end{table}

\chapter{DRIVE-TEST PRE-TRIP CHECKLIST}

\begin{enumerate}[leftmargin=*]
    \item Phone (Poco~X3) charged to $100\%$; power bank available.
    \item Network Cell Info Lite installed; logging interval $= 1$~s;
    CSV format; neighbour cells \textbf{ON}; GPS \textbf{ON}.
    \item Phone set to \textbf{LTE only}; Wi-Fi \textbf{OFF}; battery
    optimisation \textbf{disabled} for the logger.
    \item Location mode set to \textbf{High Accuracy}.
    \item One-minute test log in the campus courtyard; CSV opened on
    laptop; row count and PCI variation checked.
    \item Mounted on handlebar; screen kept \textbf{ON}; helmet on.
    \item Start logging $30$~s before moving; stop only after returning.
    \item Export CSV; rename as
    \texttt{routeN\_direction\_period\_date.csv}; back up to the
    project repository.
\end{enumerate}

\chapter{QUESTIONS FOR NTC LTE ENGINEERS}

The structured interview with NTC LTE engineers in Pokhara will cover
the following question groups. The intent is to validate the assumed
RAN configuration in the simulator, to clarify which OSS counters and
parameters might be shareable for academic use, and to understand the
current MRO / MLB and O-RAN posture of the operator.

\textbf{A. RAN and Core}
\begin{enumerate}[leftmargin=*]
    \item Which RAN vendors are currently deployed in the Pokhara
    cluster (for example, Huawei, ZTE, Nokia, Ericsson)? Is the
    deployment mixed across sites?
    \item What is the architecture of the core (flat EPC, regional
    MMEs / SGW-LBOs)?
    \item Are X2-based handovers used cluster-internally? Where do S1-
    based handovers come into play?
    \item What is the typical inter-site distance in (a) the urban
    core (Lakeside / Mahendrapul / Prithvi Chowk), (b) the Prithvi
    Highway corridor and (c) the hilly cluster
    (Bindabasini / Sarangkot)?
\end{enumerate}

\textbf{B. Handover Configuration}
\begin{enumerate}[leftmargin=*,resume]
    \item What are the currently configured values of the A3 offset
    and the Time-to-Trigger (TTT) in the Pokhara cluster? Are they
    cluster-wide or per-cell?
    \item Is the Cell Individual Offset (CIO) used per neighbour?
    Within what range?
    \item Which thresholds are used for A2 (serving becomes worse than
    threshold)?
    \item What hysteresis values are configured?
    \item Are different settings used for highway versus urban
    versus hilly clusters? How often are they re-tuned?
\end{enumerate}

\textbf{C. Mobility Robustness and Load Balancing}
\begin{enumerate}[leftmargin=*,resume]
    \item Is Mobility Robustness Optimization (MRO) enabled? Manual or
    automatic?
    \item Is Mobility Load Balancing (MLB) enabled? Manual or
    automatic? Which counters drive the decisions?
    \item What ping-pong / RLF rates are typically seen in the urban
    core during peak hours?
    \item Is sticky-cell or PCI confusion a known issue in any
    specific Pokhara sector?
    \item How often does the optimisation team re-tune the cluster?
\end{enumerate}

\textbf{D. KPIs and OSS Counters}
\begin{enumerate}[leftmargin=*,resume]
    \item Which KPIs / counters are tracked at the OSS level
    (handover success rate, ping-pong rate, RLF, PRB utilization,
    connected-UE count, throughput per cell)?
    \item What is the granularity of these counters (per cell,
    per minute, per hour)?
    \item Which of these counters could be shared, in anonymised form,
    for academic use under a clear non-disclosure scope?
\end{enumerate}

\textbf{E. Voice and Coverage}
\begin{enumerate}[leftmargin=*,resume]
    \item What is the approximate share of VoLTE versus CSFB in the
    Pokhara cluster?
    \item Where are the known coverage holes or weak-signal pockets
    (urban / highway / hilly)?
\end{enumerate}

\textbf{F. AI/ML and O-RAN Roadmap}
\begin{enumerate}[leftmargin=*,resume]
    \item Is any AI- or ML-based optimization currently in production
    in NTC's LTE network?
    \item Is NTC evaluating O-RAN-style disaggregated RAN or Near-RT
    RIC technology? What is the indicative timeline?
    \item Would NTC be interested in piloting an academic xApp on a
    test cell?
\end{enumerate}

\textbf{G. Data Sharing for the Project}
\begin{enumerate}[leftmargin=*,resume]
    \item Would NTC be willing to share, for a clearly defined academic
    window, (a) anonymised handover-event logs, (b) per-cell
    PRB-utilization counters, (c) the A3 / TTT / CIO parameter table for
    a small set of sites in Pokhara?
    \item Under what NDA / data-handling terms would such sharing be
    possible?
    \item Who in the Pokhara cluster would be the point of contact for
    the data exchange and signoff?
\end{enumerate}

\chapter{FORMAL DATA REQUEST LETTER TO NEPAL TELECOM}

\noindent\textbf{Subject:} Request for Anonymised LTE Drive-Test, OSS
and Configuration Data for an Academic Major Project.

\vspace{0.6cm}
\noindent To,\\
The Branch Manager,\\
Nepal Telecom (NTC), Pokhara Branch,\\
Kaski, Nepal.

\vspace{0.4cm}
\noindent Through:\\
The Head of Department,\\
Department of Electronics and Computer Engineering,\\
Paschimanchal Campus, Institute of Engineering,\\
Tribhuvan University, Lamachour, Pokhara--16.

\vspace{0.4cm}
\noindent Respected Sir/Madam,

We are a group of final-year students of the Bachelor of Electronics,
Communication and Information Engineering at Paschimanchal Campus,
Institute of Engineering, Tribhuvan University. As part of our major
project titled \textit{``SON-GNN-DQN: A Topology-Generalized Graph
Neural Network and Deep Q-Learning Framework for Predictive Handover
and Load Optimization in LTE Networks''}, we have already designed and
implemented an LTE mobility simulator, a Graph Neural Network
preference model, a Deep Q-Learning decision agent and a safety-bounded
Self-Organizing Network translation layer that produces bounded Cell
Individual Offset and Time-to-Trigger updates.

The next phase of our project requires validation of the framework
against measurements collected on a live LTE network. We respectfully
request Nepal Telecom's cooperation in sharing, strictly for academic
and non-commercial use, the following data for a small set of sites in
the Pokhara cluster:

\begin{enumerate}[leftmargin=*]
    \item Anonymised handover-event logs (without subscriber
    identifiers) covering a representative time window in the urban
    core, the Prithvi Highway corridor and the hilly cluster.
    \item Per-cell counters relevant to the project, including
    handover success rate, ping-pong rate, Radio Link Failure rate,
    Physical Resource Block utilization and connected-UE count, at a
    granularity that NTC considers shareable (for example, hourly
    summaries).
    \item The current Event~A3 offset, Time-to-Trigger and Cell
    Individual Offset configuration table for the chosen sites.
    \item Any guidance on permissible drive-test logging using a
    consumer Android handset on the NTC LTE network for the same
    routes.
\end{enumerate}

All data will be used exclusively for academic analysis, will be
stored in encrypted form on the project's private repository, and will
not be reproduced in any public document beyond aggregated, anonymised
statistics in our final project report and oral defence. We are happy
to sign a Non-Disclosure Agreement on terms acceptable to NTC, and to
share the final findings with the Pokhara optimisation team.

We sincerely appreciate Nepal Telecom's long-standing contribution to
academic research in Nepal and hope that this collaboration can support
both our learning objectives and NTC's continuing optimisation efforts.

\vspace{0.6cm}
\noindent Yours faithfully,

\vspace{0.6cm}
\begin{tabular}{p{6cm}p{6cm}}
Rojin Puri \dotfill   & Rujan Subedi \dotfill \\[0.5cm]
(PAS078BEI028)        & (PAS078BEI030) \\[0.6cm]
Saange Tamang \dotfill & Sulav Kandel \dotfill \\[0.5cm]
(PAS078BEI031)        & (PAS078BEI041) \\
\end{tabular}

\vspace{0.6cm}
\noindent Endorsed by:\\[0.4cm]
\noindent\rule{6cm}{0.5pt}\\
Asst.\ Prof.\ Er.\ Hom Nath Tiwari\\
Project Supervisor and Head of Department\\
Department of Electronics and Computer Engineering\\
Paschimanchal Campus, Institute of Engineering\\
Lamachour, Pokhara--16, Nepal.

\end{document}
